% ==============================================================================
% פרק 7: סיכום, תובנות ומבט לעתיד
% ==============================================================================
\chapter{סיכום, תובנות ומבט לעתיד}
\label{chap:conclusion}

\section{סינתזה אינטגרטיבית של דוח ההכנה}
\label{sec:conclusion_synthesis}
דוח מכין זה הציג תשתית תיאורטית, מתודולוגית והנדסית שלמה לפרויקט \projectname{} -- פלטפורמת מחקר עתידית מבוססת חברה רב-סוכנית אוטונומית, המיועדת להתמודד עם אחת התופעות הפסיכולוגיות המורכבות ביותר של העידן הדיגיטלי: שחיקת הרווחה הנפשית של צעירים כתוצאה מהשוואה חברתית בפידים אלגוריתמיים.

המסע המחקרי שנפרס לאורך הדוח גישר על פערים מסורתיים שבין פסיכולוגיה חברתית למדעי המחשב:
\begin{itemize}
    \item \textbf{ברמה התיאורטית}: נותח השיבוש המבני שהחילו מנועי ההמלצה המודרניים על תיאוריית ההשוואה החברתית של פסטינגר~\cite{festinger1954theory}. הודגם כיצד אופטימיזציה עסקית למקסום עוררות מציפה את הפיד באידיאליזציה קיצונית, מעצימה את אפקט הניגוד, ומרחיבה את פערי העצמי לפי היגינס~\cite{higgins1987self}.
    \item \textbf{ברמה המתודולוגית}: נחשף הכשל המבני של כלי הרווחה המסחריים הקיימים, הנשענים על מדידה כמותית עיוורת של "זמן מסך כולל" ומעוררים התנגדות פסיכולוגית (\textenglish{Reactance}). כמענה לכך, גובש מודל סימולציה מבוקר עם 20 פרסונות גנרטיביות המאפשר לבודד קשרים סיבתיים בבטיחות אתית מלאה.
    \item \textbf{ברמה ההנדסית}: פותחה ארכיטקטורה מרובת סוכנים מתואמת (\agentObserver{}, \agentAnalyst{}, \agentReflector{}, \agentInterventionist{}) הפועלת מקומית בקצה (\textenglish{Local-First Edge AI}), ממדלת את וקטור המצב הנפשי $\wellbeingVector{}$, ומספקת פיגומים מטא-קוגניטיביים עדינים המחזירים למשתמש את השליטה הפנימית.
\end{itemize}

\section{תרומה מחקרית ויישומית למדעי המחשב החברתיים ובינה מלאכותית מרובת-סוכנים}
\label{sec:conclusion_impact}
למחקר מכין זה השלכות משמעותיות הן על המחקר האקדמי והן על התעשייה היישומית:
\begin{enumerate}
    \item \textbf{קידום מדעי המחשב החברתיים (\textenglish{Computational Social Science})} -- הפרויקט מדגים כיצד מודלי שפה אוטונומיים יכולים לשמש לא רק ככלי עיבוד שפה, אלא כמצע מחקרי למידול דינמיקות רגשיות מורכבות ברזולוציה זמנית חסרת תקדים.
    \item \textbf{הגדרת פרדיגמה חדשה לרווחה דיגיטלית (\textenglish{Digital Well-Being 2.0})} -- מעבר מכלי ענישה וחסימה פטרנליסטיים לכלים מבוססי מודעות, רפלקציה עצמית ותמיכה באוטונומיה אישית בהתאם לעקרונות הפסיכולוגיה החיובית.
    \item \textbf{מודל מופת לאתיקה ופרטיות מקומית} -- הוכחת היתכנות לארכיטקטורת AI אינטימית המלווה את עולמו הרגשי של המשתמש מבלי לסחור בנתוניו או להעביר מידע אישי לשרתים מרכזיים.
\end{enumerate}

\section{דברי סיום: סוכנות אנושית, ביקורתיות ואחריות אלגוריתמית}
\label{sec:conclusion_remarks}
בעידן שבו אלגוריתמים מסחריים מתוכננים ללכוד את תשומת הלב האנושית ולנצל חולשות פסיכולוגיות אבולוציוניות, האתגר הגדול אינו ניתוק טכנולוגי גורף, אלא \textbf{החזרת הסוכנות והריבונות הקוגניטיבית למשתמש}.

פלטפורמת \projectname{} מוכיחה כי בינה מלאכותית אינה חייבת לשמש ככלי להעמקת ההתמכרות והאטימות; היא יכולה לשמש כ"מראה" צלולה ומעצימה, המלמדת צעירים להביט בפיד בעיניים ביקורתיות, לזהות אשליות דיגיטליות, ולפתח חוסן רגשי בעולם מקוון. דוח מכין זה מניח את היסודות המדעיים וההנדסיים להגשמת חזון זה, וסולל את הדרך למימוש מלא של המערכת בשלבי המחקר הבאים.
